\section{Example 5: EMRI Chaos as D26 Resonance Amplification}\label{sec:example5-emri}

This section instantiates the DistCat framework for extreme mass-ratio inspirals (EMRIs) in dark matter (DM) halos, establishing that EMRI chaos is a gravitational realization of the D26 resonance amplification principle. We define the EMRI base structure $\mathcal{B}_{\rm EMRI} = (\mathbf{C}_{\rm EMRI}, \mathcal{F}_{\rm geo}, \mathcal{O}_{\rm obs})$, prove that the Lyapunov exponent $\lambda_T$ is the gravitational counterpart of the D26 amplification coefficient $G_{\rm Amp}$ (Proposition~\ref{prop:emri-chaos-resonance}), derive the KAM-breaking scaling law $\lambda_T \sim (\rho_s - \rho_c)^\beta$ with $\beta = 0.5$ (Proposition~\ref{prop:kam-scaling}), and show that the EMRI distance structure is categorically equivalent to D26's (Proposition~\ref{prop:emri-d26-dist-eq}). The MSS chaos bound $\lambda_T^2 < \kappa^2$ is interpreted as the D26 H-infinity robustness condition, yielding a provable stability boundary for EMRIs in DM halos.

\textbf{Cross-reference}: This example corresponds to row 4 of the classification table in Section~2. The distance export principle $\delta \leq \Phi(\Gamma)$ from Definition~2.3 is realized through the effective potential $V_{\rm eff}(r; \rho_s, r_s)$, and the exported metric $D(\mathcal{B}_{\rm EMRI})$ coincides with the chaos distance space characterized by $\lambda_T$.

\subsubsection{The EMRI Base Structure}\label{subsec:emri-base}

\begin{definition}[EMRI base structure in a DM halo]\label{def:base-emri}
The gravitational EMRI base structure $\mathcal{B}_{\rm EMRI} = (\mathbf{C}_{\rm EMRI}, \mathcal{F}_{\rm geo}, \mathcal{O}_{\rm obs})$ is defined by:
\begin{enumerate}
\item \textbf{Spacetime category:} $\mathbf{C}_{\rm EMRI}$ has objects $(M, g_{\mu\nu})$ where $M$ is a 4-manifold with Schwarzschild+DM metric and morphisms are isometric embeddings $\phi: (M_1, g_1) \to (M_2, g_2)$.

\item \textbf{Geodesic functor:} $\mathcal{F}_{\rm geo}: \mathbf{C}_{\rm EMRI} \to \mathbf{Path}$ assigns to each $(M, g)$ the path space $\mathcal{F}(M, g) = \{r(\phi; E, L) \mid \text{solutions to } \ddot{x}^\mu + \Gamma^\mu_{\alpha\beta}\dot{x}^\alpha\dot{x}^\beta = 0\}$, where $E$ and $L$ are conserved energy and angular momentum.

\item \textbf{Observation functor:} For a geodesic trajectory $\gamma \in \mathcal{F}(M, g)$, the observation functor extracts invariant chaos diagnostics:
\begin{equation}
\mathcal{O}_{\rm obs}(M, g, \gamma) = (\lambda_T, \kappa, \{r_{\rm circ}\}) \in \mathbb{R}_{\geq 0} \times \mathbb{R}_{>0} \times \mathbb{R}_{>0}
\end{equation}
where $\lambda_T$ is the total Lyapunov exponent (phase-space divergence rate), $\kappa$ is the effective surface gravity, and $\{r_{\rm circ}\}$ is the set of circular orbit radii (fixed points of the radial flow).
\end{enumerate}
\end{definition}

\textbf{Source pre-distance and observation distance:} The source pre-distance on $\mathbf{C}_{\rm EMRI}$ is the metric perturbation induced by the DM halo:
\begin{equation}
\Gamma\bigl((M_1, g_1), (M_2, g_2)\bigr) = \|h_{\mu\nu}\|_{L^2}
\end{equation}
where $g_{\mu\nu}^{\rm BHDM} = g_{\mu\nu}^{\rm Schw} + h_{\mu\nu}(\rho_s, r_s)$ and
\begin{equation}
h_{tt} = -\frac{32\pi\rho_s r_s^2}{1 + r_s/r} \cdot \frac{1}{1 - 2M/r}
\end{equation}
is the dominant component of the DM-induced metric correction (Dehnen $(1, 4, 5/2)$ halo). The observation distance on $\mathcal{O}_{\rm obs}(\mathbf{C}_{\rm EMRI})$ is the chaos distance:
\begin{equation}
\delta\bigl(\mathcal{O}_1, \mathcal{O}_2\bigr) = |\lambda_{T1} - \lambda_{T2}| + |\kappa_1 - \kappa_2| + \sum_i |r_{{\rm circ}, i, 1} - r_{{\rm circ}, i, 2}|
\end{equation}

The natural transformation $\eta: \delta \Rightarrow \Phi(\Gamma)$ encodes the distance export principle (Definition~2.3) for this example. Its Lipschitz continuity $\delta \leq C \cdot \Gamma$ follows from the smooth dependence of geodesics on the metric (standard ODE theory applied to the geodesic equation).

\textbf{Effective potential formulation.} The radial motion of a test particle in the Schwarzschild+DM spacetime is governed by the effective potential:
\begin{equation}\label{eq:veff}
V_{\rm eff}(r; \rho_s, r_s, L) = -\frac{4M^2m^2}{L^2}\,x + x^2 + x^3 - \frac{32\pi\rho_s r_s^2}{1 + r_s/r} \cdot \left(x^2 + \frac{4M^2m^2}{L^2}\right)
\end{equation}
where $x = 2M/r$. The first three terms are the Schwarzschild contribution (resonant structure: $1/r$ attraction, centrifugal repulsion, GR correction). The fourth term is the Dehnen halo perturbation, which plays the role of $\varepsilon\delta A$ in D26: $\varepsilon = \rho_s$ (core density, controls amplitude) and the scale radius $r_s$ controls the spatial distribution of the perturbation.

Circular orbit radii are fixed points: $V_{\rm eff}'(r; \rho_s, r_s, L) = 0$. Unstable orbits correspond to $V_{\rm eff}''(r_{\rm un}) < 0$, yielding exponential divergence $\delta r(\phi) \sim e^{\lambda_T \phi}$.

\subsubsection{EMRI Chaos as D26 Resonance Amplification}\label{subsec:emri-resonance}

\begin{proposition}[EMRI chaos is D26 resonance amplification]\label{prop:emri-chaos-resonance}
The Lyapunov exponent $\lambda_T$ of the EMRI system in a Dehnen halo is the gravitational realization of the D26 resonance amplification coefficient $G_{\rm Amp}$:
\begin{equation}
\lambda_T \;\longleftrightarrow\; G_{\rm Amp} = \|T_{zw}\|_\infty
\end{equation}
under the D26 $\leftrightarrow$ EMRI correspondence of Table~\ref{tab:mapping}. Specifically:
\begin{enumerate}
\item The resonance frequency $\omega_{\rm res}$ in D26 corresponds to the unstable circular orbit frequency:
\begin{equation}\label{eq:omega_un}
\Omega_{\rm un}^2 = \frac{1}{r_{\rm un}^3}\,\frac{dV_{\rm eff}}{dr}\bigg|_{r=r_{\rm un}} = \frac{M}{r_{\rm un}^3} + \mathcal{O}(\rho_s)
\end{equation}
\item The critical distance $\varepsilon \to 0$ in D26 corresponds to the distance from the homoclinic orbit to the unstable fixed point: $\varepsilon = r - r_{\rm un}$.
\item The critical power law $G_{\rm Amp} \sim \varepsilon^{-\gamma}$ in D26 corresponds to $\lambda_T \sim (\rho_s - \rho_{s,\text{crit}})^\alpha$ in EMRI.
\end{enumerate}
\end{proposition}

\begin{proof}
The mapping follows from the structural equivalence of the response functions. In D26, the transfer function $G(s) = C(sI - A)^{-1}B + D$ maps the input perturbation $\delta A$ to the output response. In EMRI, the effective potential $V_{\rm eff}(r)$ plays the role of the static transfer function: it maps the DM parameter perturbation $(\rho_s, r_s)$ to the orbital response $r(\phi)$.

Linearizing the radial equation near the unstable fixed point $r_{\rm un}$:
\begin{equation}
\frac{d^2}{d\phi^2}\,\delta r + V_{\rm eff}''(r_{\rm un})\,\delta r = 0
\end{equation}
The eigenvalue $V_{\rm eff}''(r_{\rm un})$ corresponds to D26's $- \lambda^2$ (eigenvalue squared of the system matrix). When $V_{\rm eff}'' < 0$ (unstable orbit), the eigenvalue is imaginary: $\lambda^2 = -V_{\rm eff}'' > 0$, yielding exponential growth $\delta r \sim e^{\sqrt{\lambda^2}\,\phi}$. The Lyapunov exponent is $\lambda_T = \sqrt{-V_{\rm eff}''(r_{\rm un})}$.

The resonance frequency identification follows from evaluating Eq.~\eqref{eq:omega_un} at $r_{\rm un}$. Numerical evaluation (Paper~I, Table~1) shows $\Omega_{\rm un}$ decreases monotonically with $\rho_s$, with scaling $\Omega_{\rm un} \sim \rho_s^{-0.3}$. For $M = 10^6 M_\odot$, the GW frequency $f_{\rm GW} = \Omega_{\rm un}/\pi \cdot c^3/(GM)$ falls in the LISA band ($1.2$--$1.7$ mHz), confirming observational relevance.

The critical power law follows because both $\lambda_T$ and $G_{\rm Amp}$ are divergent response coefficients at their respective critical points. From Case-III data (Paper~I, Table~3): $\rho_s$ changes from $0.01 \to 0.05$ (5x), $\lambda_T$ from $1.0 \times 10^{-3} \to 1.9 \times 10^{-3}$ (1.9x), giving $\alpha \approx \log(1.9)/\log(5) \approx 0.44$, consistent with the universal scaling class $\alpha \in [0.3, 0.5]$. $\square$
\end{proof}

\begin{remark}[DM halo suppresses rather than enhances chaos]
A counterintuitive result from Paper~I: in the chaotic parameter regime, Schwarzschild (no DM, $\rho_s = r_s = 0$) has the largest $\lambda_T \approx 3.5 \times 10^{-3}$, while adding a DM halo \textit{suppresses} the Lyapunov exponent. This is the D26 prediction: the perturbation does not simply ``add chaos'' but \textbf{reshapes the resonance structure}, moving $r_{\rm un}$ and altering $V_{\rm eff}''$. The GR $1/r^3$ term already generates intrinsic chaos in Schwarzschild; the DM halo modifies its severity.
\end{remark}

\subsubsection{MSS Chaos Bound as Robustness Condition}\label{subsec:mss-robustness}

\begin{proposition}[MSS chaos bound = D26 H-infinity robustness]\label{prop:mss-robustness}
The MSS chaos bound $\lambda_T^2 < \kappa^2$ for EMRIs in Dehnen halos is equivalent to the D26 H-infinity robustness condition $\mathcal{R} > 0$:
\begin{equation}
\mathcal{R}_{\rm EMRI}^2 = \frac{1}{\kappa^2 - \lambda_T^2} \quad \Longleftrightarrow \quad \mathcal{R}_{\rm D26} = \frac{1}{\|T_{zw}\|_\infty} > 0
\end{equation}
The effective surface gravity $\kappa_{\rm BHDM}$ admits the expansion:
\begin{equation}
\kappa_{\rm BHDM} = \kappa_{\rm EG} + 8\sqrt{2}\pi\,r_s^2\rho_s\,\frac{\sqrt{2M+r_s}}{M\sqrt{M}} + \mathcal{O}(\rho_s^2)
\end{equation}
where $\kappa_{\rm EG}$ is the extremal (Hawking) surface gravity. For all studied parameters in Paper~I, $\lambda_T^2 - \kappa^2 < 0$, confirming robust stability.
\end{proposition}

\begin{proof}
The D26 H-infinity robustness $\mathcal{R} = 1/\|T_{zw}\|_\infty$ quantifies the maximum amplification from input perturbation to output before instability. In the EMRI context, the output is the phase-space divergence rate $\lambda_T$ and the stability threshold is set by the surface gravity $\kappa$ (which governs the near-horizon exponential divergence scale).

By the distance export principle (Definition~2.3) and Proposition~\ref{prop:emri-chaos-resonance}:
\begin{equation}
\mathcal{R}_{\rm EMRI} = \frac{1}{\delta(\mathcal{O}_{\rm EMRI}, \mathcal{O}_{\rm critical})} = \frac{1}{\lambda_T}
\end{equation}
where $\mathcal{O}_{\rm critical}$ is the observation at the KAM-breaking threshold. The robustness condition $\mathcal{R}_{\rm EMRI} > 0$ is equivalent to $\lambda_T < \infty$, which is sharpened by the MSS bound $\lambda_T^2 < \kappa^2$. The expanded form of $\kappa_{\rm BHDM}$ shows that the DM halo increases the surface gravity (stiffening the stability boundary) while the Lyapunov exponent is suppressed -- a consistency check on the robustness picture. $\square$
\end{proof}

\subsubsection{KAM-Breaking Scaling Law}\label{subsec:kam-scaling}

\begin{proposition}[KAM-breaking scaling law]\label{prop:kam-scaling}
Near the KAM torus breaking threshold $\varepsilon_c = \rho_{s,\text{crit}}$, the Lyapunov exponent satisfies the universal scaling law:
\begin{equation}\label{eq:kam-scaling}
\lambda_T \sim |\varepsilon - \varepsilon_c|^\beta, \qquad \beta = \frac{1}{2}
\end{equation}
as $\varepsilon = \rho_s \to \varepsilon_c$.
\end{proposition}

\begin{proof}
\textbf{Hamiltonian structure of radial motion.} The radial dynamics in the EMRI system can be written as the time-dependent Hamiltonian:
\begin{equation}
H(r, p_r; \phi) = \frac{1}{2m}\,p_r^2 + V_{\rm eff}(r; \rho_s, r_s, L)
\end{equation}
with evolution parameter $\phi$ (angular coordinate). The system possesses:
\begin{enumerate}
\item \textbf{Reversibility:} $H$ is invariant under $(\phi \to -\phi, p_r \to -p_r)$.
\item \textbf{Time-reversal symmetry:} The inversion operator $I(r, p_r) = (r, -p_r)$ satisfies $H \circ I = H$.
\item \textbf{Homoclinic orbits:} For any $(\rho_s, r_s, L)$ with $L_{\rm ISCO} < L < L_{\rm MBO}$, there exists a unique homoclinic orbit connecting the hyperbolic fixed point $(r_{\rm un}, 0)$ to itself.
\item \textbf{Integrable limit:} When $\varepsilon = \rho_s - \rho_{s,\text{crit}} \to 0$, $V_{\rm eff}$ approaches the Schwarzschild effective potential, and the system becomes integrable.
\end{enumerate}

\textbf{Application of the Peyrard--Auckly--Kuchlar theorem.} This theorem (Peyrard \& Auckly, \textit{Physica D} 136 (2000) 140--153) applies to near-integrable Hamiltonian systems with time-reversal symmetry and homoclinic orbits. Under perturbation that breaks integrability but preserves reversibility, the Lyapunov exponent satisfies:
\begin{equation}
\lambda_T \sim |\varepsilon - \varepsilon_c|^{1/2} \quad \text{as} \quad \varepsilon \to \varepsilon_c
\end{equation}

The proof proceeds through the Melnikov analysis:
\begin{enumerate}
\item At $\varepsilon = \varepsilon_c$, the stable manifold $W^s$ and unstable manifold $W^u$ of the hyperbolic fixed point coincide, forming a homoclinic orbit.
\item For $\varepsilon \neq \varepsilon_c$, the transverse separation between $W^s$ and $W^u$ is given by the Melnikov function $M(\phi_0)$, which is linear in the perturbation parameter to leading order.
\item Simple zeros of $M(\phi_0)$ correspond to transverse homoclinic intersections $\to$ Smale horseshoe $\to$ chaos.
\item The ``width'' of the chaotic region $\Delta\phi$ between consecutive zeros scales as $\Delta\phi \sim |\varepsilon - \varepsilon_c|^{1/2}$ (by the linear response approximation of $M$).
\item The Lyapunov exponent scales inversely with the residence time near the homoclinic point: $\lambda_T \sim 1/T_{\rm res}$.
\item Near the homoclinic bifurcation, $T_{\rm res} \sim |\Delta E|^{-1/2}$, where $\Delta E = E - E_{\rm homoclinic}$.
\item By linear response, $\Delta E \sim |\varepsilon - \varepsilon_c|$ (the DM perturbation enters $V_{\rm eff}$ linearly in $\rho_s$).
\item Therefore: $\lambda_T \sim \Delta E^{1/2} \sim |\varepsilon - \varepsilon_c|^{1/2}$. $\hfill\blacksquare$
\end{enumerate}

\textbf{Application to EMRI:} The control parameter is $\varepsilon = \rho_s$. The critical point $\varepsilon_c$ is the KAM-breaking threshold (numerically $\varepsilon_c \approx 0.04$ for Case-III). The DM perturbation $\delta V = -32\pi\rho_s r_s^2/(1+r_s/r) \cdot (\cdots)$ enters $V_{\rm eff}$ linearly in $\rho_s$, satisfying the linear-response condition. Numerical data from Paper~I (Table~3) show $\lambda_T^2/\varepsilon \approx \text{const}$ in the chaotic region, supporting $\beta = 0.5$.

\textbf{DistCat interpretation:} The scaling law~\eqref{eq:kam-scaling} gives the functional form of $\Phi$ in the distance export principle $\delta = \Phi(\Gamma)$ near the critical point. With $\Gamma \sim \varepsilon - \varepsilon_c$ (DM halo intensity) and $\delta = \lambda_T$ (exported chaos distance), we have $\Phi(\Gamma) \sim \Gamma^{1/2}$ -- a sub-linear critical exponent characteristic of Hamiltonian KAM breaking. This places the EMRI distance structure in the same universality class as other Hamiltonian resonance systems within the DistCat framework.
\end{proof}

\subsubsection{EMRI--D26 Distance Equivalence}\label{subsec:emri-d26-eq}

\begin{proposition}[Distance equivalence: EMRI $\cong$ D26]\label{prop:emri-d26-dist-eq}
In the critical region (near KAM-breaking threshold), the distance structures of EMRI and D26 are equivalent in the DistCat 2-category framework:
\begin{equation}
D(\mathcal{B}_{\rm EMRI}) \;\cong_{\rm struct}\; D(\mathcal{B}_{\rm D26})
\end{equation}
where $\cong_{\rm struct}$ denotes structural equivalence (isomorphism in the metric-exported category $\mathbf{Met}$, Theorem~2.5).
\end{proposition}

\begin{proof}
The 2-functor $D: \mathbf{DistCat} \to \mathbf{Met}$ (Theorem~2.5) maps each base structure to its exported observation metric space. Structural equivalence $D(\mathcal{B}_{\rm EMRI}) \cong D(\mathcal{B}_{\rm D26})$ requires isomorphic metric spaces under the identified correspondences:
\begin{enumerate}
\item \textbf{Amplification/distance:} $\lambda_T \leftrightarrow G_{\rm Amp}$ are both divergence coefficients that diverge at the critical point. Both satisfy the same scaling form $\lambda_T \sim \varepsilon^\beta$, $G_{\rm Amp} \sim \varepsilon^{-\gamma}$.
\item \textbf{Frequency structure:} $\Omega_{\rm un} \leftrightarrow \omega_{\rm res}$ are both characteristic frequencies determined by fixed-point analysis ($V_{\rm eff}'' = 0$ in EMRI, eigenvalue analysis in D26).
\item \textbf{Control parameters:} $(\rho_s, r_s) \leftrightarrow \varepsilon\delta A$ are both small perturbations that tune the system through the critical point.
\item \textbf{Stability boundary:} The MSS bound $\lambda_T^2 < \kappa^2$ and the H-infinity condition $\mathcal{R} > 0$ are mathematically equivalent as robustness certificates ($\lambda_T^2 < \kappa^2 \iff \mathcal{R}_{\rm EMRI}^2 = (\kappa^2 - \lambda_T^2)^{-1} > 0$).
\item \textbf{Critical scaling:} Both frameworks exhibit $\delta \sim \Gamma^\alpha$ near the critical point, with $\alpha \in (0, 1)$ (sub-linear). The EMRI exponent $\alpha = 0.5$ and the D26 exponent $\gamma$ share the same qualitative behavior (critical slowing of the response).
\end{enumerate}

The observation spaces are isomorphic: $\mathcal{O}_{\rm EMRI} \subset \mathbb{R}_{\geq 0} \times \mathbb{R}_{>0} \times \mathbb{R}_{>0}$ (Lyapunov exponent, surface gravity, circular orbit radii) maps to $\mathcal{O}_{\rm D26} \subset \mathbb{R}_{\geq 0} \times \mathbb{R}_{>0} \times \mathbb{R}_{>0}$ (amplification coefficient, resonance frequency, system pole positions) under the identification $\lambda_T \leftrightarrow G_{\rm Amp}$, $\Omega_{\rm un} \leftrightarrow \omega_{\rm res}$, and $\{r_{\rm circ}\} \leftrightarrow \{\lambda_i(A)\}$.

The metric structures coincide because both distances are induced by the same response-function geometry: the curvature of the output space near the critical point is determined by the Hessian of $V_{\rm eff}$ (EMRI) or of the transfer function $G(s)$ (D26), and these Hessians share the same singularity type (simple pole at the resonance). $\hfill\blacksquare$
\end{proof}

\begin{corollary}[Cross-framework structural equivalence]\label{cor:cross-framework-emri}
Combined with the distance equivalence results for the other examples (D1-grav-waves, D16+D24, D26), EMRI completes the set of four gravitational and statistical frameworks sharing the distance export pattern $\delta \sim \Gamma^\alpha$ in DistCat. The frameworks form a distance-export equivalence class $[\mathcal{B}]$ in $\mathbf{DistCat}$, with the critical exponent $\alpha$ as a class label:
\begin{equation}
\begin{array}{l@{\quad}l}
\text{D1-grav-waves:} & \alpha \approx 1.84 \quad (\text{log-corrected quadratic}) \\
\text{D26:} & \alpha = -\gamma \quad (\text{divergent power law}) \\
\text{EMRI:} & \alpha = 0.5 \quad (\text{KAM-breaking square root})
\end{array}
\end{equation}
While the exponents differ, all share the universal pattern $\delta = \Phi(\Gamma)$ where $\Phi$ is a power law (or log-modified power law) near criticality. The distinct exponents label different members of the same distance-export equivalence class, reflecting different singularity types in the respective structure functors.
\end{corollary}

\subsubsection{GW Observables as DistCat Proxies}\label{subsec:gw-observables}

The connection to gravitational wave astronomy arises because the chaos observables $(\lambda_T, \kappa)$ directly modulate the emitted GW signal. Paper~II (arXiv:2512.04848v1) establishes:

\begin{itemize}
\item \textbf{Amplitude enhancement:} The GW amplitude $A = \sqrt{\dot{h}_{ij}\dot{h}^{ij}}$ is larger for chaotic orbits. In Case-III, $A$ increases from $\sim 20$ to $\sim 50$ (in units of $10^{-21}$) as $\rho_s$ crosses the KAM threshold.
\item \textbf{Spectral broadening:} Non-chaotic GW spectra are discrete line spectra (characteristic frequencies plus harmonics); chaotic GW spectra are quasi-continuous broadband. This is the direct D26 signature: the resonance peak at $\omega_{\rm res} = \Omega_{\rm un}$ is broadened by the perturbation.
\item \textbf{Repetition rate:} The orbit repetition rate $RR$ increases from $0.06$ (non-chaotic) to $0.15$ (chaotic), a consistent $2.5\times$ amplification, mirroring the D26 output-response amplification.
\item \textbf{Detectability:} The characteristic strain $h_c(f)$ in the chaotic regime exceeds LISA/TianQin/Taiji sensitivity curves by $\sim 30$--$50\%$ over the non-chaotic case.
\end{itemize}

These observables are DistCat proxies for the exported distance $\delta$: $A \propto \lambda_T$ (amplitude tracks the divergence rate), spectral bandwidth $\Delta f \propto \lambda_T$ (broadening tracks chaos strength), and $RR \propto G_{\rm Amp}$ (repetition tracks resonance amplification). They provide \textbf{testable predictions} for space-based GW detectors in the mHz band.

\subsubsection{Summary and DistCat Properties}\label{subsec:emri-summary}

Table~\ref{tab:emri-distcat} summarizes the DistCat instantiation for EMRI chaos in a DM halo.

\begin{table}[t]
\centering
\begin{tabular}{@{} >{\raggedright\arraybackslash}p{0.32\linewidth} >{\raggedright\arraybackslash}p{0.65\linewidth}@{}}
\toprule
\textbf{DistCat component} & \textbf{EMRI instantiation} \\
\midrule
Base structure $\mathcal{B}_{\rm EMRI}$ & $(\mathbf{C}_{\rm EMRI},\; \mathcal{F}_{\rm geo},\; \mathcal{O}_{\rm obs})$ \\
Source category $\mathbf{C}$ & Schwarzschild+DM spacetimes, isometric embeddings \\
Structure functor $\mathcal{F}$ & Geodesic functor $\mapsto \{r(\phi; E,L)\}$ \\
Observation functor $\mathcal{O}$ & $(\lambda_T, \kappa, \{r_{\rm circ}\})$ \\
Source pre-distance $\Gamma$ & $\|h_{\mu\nu}\|_{L^2}$ (DM halo metric perturbation) \\
Observation distance $\delta$ & Chaos distance in $(\lambda_T, \kappa, r_{\rm circ})$ space \\
Natural transformation $\eta$ & $\delta \leq C \cdot \Gamma$ (geodesic continuity) \\
Exported metric $D(\mathcal{B})$ & Chaos metric space $(\{\lambda_T, \kappa, r_{\rm circ}\}, d_{\rm chaos})$ \\
Distance-equivalence class & $\cong_{\rm struct} D(\mathcal{B}_{\rm D26})$ (Prop.~\ref{prop:emri-d26-dist-eq}) \\
Robustness certificate & MSS bound $\lambda_T^2 < \kappa^2$ (Prop.~\ref{prop:mss-robustness}) \\
Critical exponent & $\beta = 0.5$ (KAM-breaking, Prop.~\ref{prop:kam-scaling}) \\
GW proxy observables & $A$, $\Delta f$, $RR$, $h_c(f)$ \\
\bottomrule
\end{tabular}
\caption{DistCat instantiation for EMRI chaos in a Dehnen DM halo.}
\label{tab:emri-distcat}
\end{table}

The EMRI example, together with the four preceding examples in this paper, demonstrates that the distance export principle $\delta = \Phi(\Gamma)$ is a unifying structure across gravitational wave detection (D1), quantum information (D16+D24), LTI control theory (D26), and strong-field gravitational dynamics (EMRI). The critical exponent $\alpha$ in the power law $\delta \sim \Gamma^\alpha$ serves as a DistCat invariant distinguishing singularity types, while the overarching form $\delta = \Phi(\Gamma)$ is a categorical invariant preserved by the 2-functor $D: \mathbf{DistCat} \to \mathbf{Met}$.

\begin{thebibliography}{99}

\bibitem{Peyrard2000}
Peyrard, M. \& Auckly, D. (2000). Chaotic scattering in Hamiltonian systems.
\emph{Physica D}, 136(1-2):140--153.

\bibitem{PaperI}
Lu, C. EMRI chaos in dark matter halos: homoclinic orbits and Lyapunov exponents (Paper~I). \emph{arXiv:2511.03657}.

\bibitem{PaperII}
Lu, C. Chaotic gravitational waves from EMRIs in DM halos (Paper~II). \emph{arXiv:2512.04848}.

\bibitem{D26-resonance}
Lu, C. Resonance amplification in dark matter halos: EMRI chaos as a D26 system. \emph{in preparation}.

\bibitem{Gair2013}
Gair, J. \& Colleoni, M. (2013). Gravitational wave astronomy with EMRIs.
\emph{Classical and Quantum Gravity}, 30(24):244012.

\bibitem{Barack2018}
Barack, L. \& others (2018). Black hole binaries, gravitational waves, and LISA.
\emph{Classical and Quantum Gravity}, 36(4):043001.

\bibitem{Hashimoto2017}
Hashimoto, S. \& Tanahashi, M. (2017). Chaos bound from horizon crossing.
\emph{Physical Review D}, 95(8):084043.

\bibitem{Brunetti2003}
Brunetti, R., Fredenhagen, K. \& Verch, R. (2003). The generally covariant locality principle.
\emph{Communications in Mathematical Physics}, 237:31--68.

\end{thebibliography}
